\documentclass[12pt,a4paper]{article}
\usepackage[utf8]{inputenc}
\usepackage[T1]{fontenc}
\usepackage[indonesian]{babel}
\usepackage{geometry}
\usepackage{longtable}
\usepackage{hyperref}
\usepackage{enumitem}
\usepackage{booktabs}
\usepackage{float}

\geometry{margin=2.5cm}

\title{Laporan Desain Sistem: MediChart EMR v2}
\author{Kelompok Proyek MediChart}
\date{\today}

\begin{document}

\maketitle

\section*{Ringkasan Eksekutif}
\addcontentsline{toc}{section}{Ringkasan Eksekutif}
MediChart EMR v2 adalah prototipe sistem rekam medis elektronik berbasis browser (Single-Page Application) yang dirancang untuk tujuan edukasi. Dokumen ini menjelaskan desain sistem secara lengkap: tujuan, arsitektur, skema data, kebijakan keamanan berbasis peran (RBAC), modul diagnostik, serta rincian implementasi teknis yang memudahkan pengujian dan presentasi.

\tableofcontents
\newpage

\section{Pendahuluan}
MediChart EMR v2 dikembangkan sebagai proyek kelompok untuk mempraktikkan desain sistem informasi kesehatan. Fokusnya adalah menyajikan alur kerja klinis sederhana (manajemen pasien, encounter, catatan SOAP) dan menambah modul pencatatan hasil diagnostik (laboratorium dan radiologi) dengan kontrol akses yang eksplisit.

\section{Deskripsi Sistem dan Ruang Lingkup}
Sistem berjalan sepenuhnya di browser tanpa backend. Ruang lingkup fungsional meliputi:
\begin{enumerate}[leftmargin=1.2cm]
  \item Autentikasi dan manajemen sesi.
  \item Role-Based Access Control (RBAC) untuk peran: Dokter, Perawat, Admin.
  \item CRUD data pasien dan encounter.
  \item Catatan klinis (SOAP) dan Problem List.
  \item Diagnostic Reports: input lab/imaging, penandaan status, interpretasi.
  \item Manajemen pengguna (Admin).
\end{enumerate}

Catatan: desain ditujukan untuk demonstrasi akademik; bukan produksi klinis.

\section{Arsitektur Sistem}
Arsitektur utama adalah SPA monolitik (semua kode berada di \texttt{medichart.html}). Struktur logical:
\begin{itemize}[leftmargin=1.2cm]
  \item Presentasi: LoginPage, AppShell, Panel (Patients, DiagReports, Admin)
  \item Logika: Auth\_Module, Patient\_Module, Clinical\_Module, DiagReport\_Module, Admin\_Module
  \item Data: localStorage (data permanen), sessionStorage (sesi aktif)
\end{itemize}

Keterangan teknis singkat:
\begin{itemize}[leftmargin=1.2cm]
  \item Keys di localStorage: \texttt{mc_patients}, \texttt{mc_encounters}, \texttt{mc_clinical_notes}, \texttt{mc_problems}, \texttt{mc_diag_reports}, \texttt{mc_users}.
  \item Sesi disimpan pada \texttt{sessionStorage['mc_session']} dengan objek `{ userId, username, name, role, loginAt }`.
\end{itemize}

\section{Model Data}
Deskripsi singkat entitas penting dan contoh struktur JSON.
\subsection*{DiagnosticReport (ringkasan)}
\begin{verbatim}
{
  "reportId": "uuid",
  "patientId": "uuid",
  "encounterId": "uuid"|null,
  "type": "lab"|"imaging",
  "date": "YYYY-MM-DD",
  "status": "pending"|"final"|"amended",
  "interpretation": "string",
  "labData": [{ "paramName":"Hb", "value":"13.5", "unit":"g/dL", "referenceRange":"12-16" }],
  "imagingData": null
}
\end{verbatim}

Relasi utama: setiap DiagnosticReport memiliki foreign key ke Patient dan opsional ke Encounter.

\section{Keamanan dan RBAC}
Kebijakan akses dirumuskan untuk tiap peran. Tabel ringkas hak akses:
\begin{table}[H]
\centering
\begin{tabular}{lccc}
\toprule
Fitur & Dokter & Perawat & Admin \\
\midrule
Lihat pasien & \checkmark & \checkmark & \checkmark \\
Tambah/Edit pasien & \checkmark &  & \checkmark \\
Input SOAP (Subjective/Assessment/Plan) & \checkmark &  &  \\
Input SOAP (Objective/Vitals) & \checkmark & \checkmark &  \\
Lihat Diagnostic Reports & \checkmark & \checkmark &  \\
Tambah/Edit Diagnostic Reports & \checkmark &  &  \\
Kelola pengguna &  &  & \checkmark \\
\bottomrule
\end{tabular}
\caption{Matriks RBAC (ringkasan)}
\end{table}

Implementasi: fungsi \texttt{canAccess(feature)} mengembalikan boolean berdasarkan role saat ini.

\section{Desain Modul dan Antarmuka (API Lokal)}
Dokumen ini juga mencantumkan nama fungsi penting yang diimplementasikan di \texttt{medichart.html} untuk memudahkan pembaca dan penguji.
\begin{itemize}[leftmargin=1.2cm]
  \item Auth\_Module: \texttt{login(username,password)}, \texttt{logout()}, \texttt{getSession()}, \texttt{hasRole(role)}, \texttt{canAccess(feature)}.
  \item Patient\_Module: \texttt{getAllPatients()}, \texttt{getPatientById(id)}, \texttt{addPatient(data)}, \texttt{updatePatient(id,data)}.
  \item Clinical\_Module: \texttt{getEncountersByPatient(pid)}, \texttt{addEncounter(data)}, \texttt{saveClinicalNote(encounterId,data)}.
  \item DiagReport\_Module: \texttt{getReportsByPatient(pid)}, \texttt{addReport(data)}, \texttt{updateStatus(reportId,status)}, \texttt{filterByType(pid,type)}, \texttt{validateReport(data)}.
  \item Admin\_Module: \texttt{getAllUsers()}, \texttt{addUser(data)}, \texttt{deactivateUser(id)}.
\end{itemize}

\section{Implementasi Teknis dan Petunjuk Pengujian}
Petunjuk singkat untuk penguji/dosen:
\begin{enumerate}[leftmargin=1.2cm]
  \item Buka file \texttt{medichart.html} di browser modern (Chrome/Firefox/Edge).
  \item Gunakan akun demo yang tersedia (lihat README) untuk login sebagai Dokter/Perawat/Admin.
  \item Uji alur: tambah pasien → buat encounter → isi SOAP → tambah Diagnostic Report (lab/imaging) → ubah status.
  \item Cek localStorage pada DevTools untuk memverifikasi koleksi: \texttt{mc_diag_reports}, \texttt{mc_patients}, dll.
\end{enumerate}

Validasi otomatis dasar: script pengecekan keberadaan file dokumentasi dan key lokal telah disiapkan oleh tim (lihat skrip validasi di repository root jika ada).

\section{Validasi dan Pengujian}
Rincian pengujian yang dilakukan:
\begin{itemize}[leftmargin=1.2cm]
  \item Smoke test login dan navigasi untuk ketiga peran.
  \item Pembuatan laporan diagnostik dengan field wajib kosong → validasi error spesifik.
  \item Transisi status: mengubah dari \texttt{pending} ke \texttt{final}; saat edit setelah \texttt{final} otomatis menjadi \texttt{amended}.
  \item Konsistensi data: jumlah yang ditampilkan sama dengan jumlah di \texttt{mc_diag_reports} untuk pasien.
\end{itemize}

\section{Petunjuk Penyertaan Diagram}
Untuk menyertakan diagram PlantUML dalam laporan PDF:
\begin{enumerate}[leftmargin=1.2cm]
  \item Render file \texttt{*.puml} di \href{https://www.plantuml.com/}{PlantUML online} atau gunakan PlantUML CLI untuk menghasilkan PNG/SVG.
  \item Simpan hasil ke folder \texttt{docs/diagrams/} (mis. \texttt{docs/diagrams/class-diagram.png}).
  \item Sisipkan pada dokumen LaTeX dengan \texttt{\includegraphics[width=\linewidth]{docs/diagrams/class-diagram.png}}.
\end{enumerate}

\section{Pembahasan dan Rekomendasi}
Ringkasan pembahasan: desain ini cocok untuk tujuan pembelajaran. Untuk langkah lanjutan (opsional) direkomendasikan:
\begin{itemize}[leftmargin=1.2cm]
  \item Migrasi data ke backend sederhana (REST API) untuk multi-user dan backup.
  \item Menambahkan enkripsi dan audit trail untuk data sensitif.
  \item Menggunakan generator diagram otomatis dalam alur build dokumen (CI) untuk menjaga konsistensi lampiran.
\end{itemize}

\section{Kesimpulan}
Laporan ini menjelaskan desain dan implementasi MediChart EMR v2 secara detil untuk kebutuhan akademik dan evaluasi. Dokumen ini dilengkapi dengan petunjuk pengujian dan cara menyertakan diagram sehingga tim dapat menyerahkan naskah dan slide presentasi.

\section*{Lampiran}
\addcontentsline{toc}{section}{Lampiran}
\begin{itemize}[leftmargin=1.2cm]
  \item \texttt{docs/makalah.md} dan \texttt{docs/makalah.tex} — naskah akademik kelompok
  \item \texttt{docs/presentation.tex} — sumber slide Beamer
  \item PlantUML sources di \texttt{docs/} (lihat README untuk rincian)
\end{itemize}

\end{document}